% ============================================================
% NOVA Sovereign Paper 11
% Capability-Driven Intelligence Certification:
% A Formal Framework for Autonomous AI Validation and Revocable Trust
% ============================================================
\documentclass[11pt,a4paper]{article}
\usepackage{amsmath,amssymb,amsthm}
\usepackage{geometry}
\usepackage{hyperref}
\usepackage{booktabs}
\usepackage{enumitem}
\usepackage{algorithm}
\usepackage{algpseudocode}

\geometry{margin=1in}

\newtheorem{theorem}{Theorem}
\newtheorem{lemma}[theorem]{Lemma}
\newtheorem{proposition}[theorem]{Proposition}
\newtheorem{corollary}[theorem]{Corollary}
\newtheorem{definition}{Definition}
\newtheorem{hypothesis}{Hypothesis}
\newtheorem{remark}{Remark}

\title{Capability-Driven Intelligence Certification:\\
A Formal Framework for Autonomous AI Validation and Revocable Trust}

\author{
  NOVA Sovereign Research Division \\
  \texttt{research@nova-sovereign.org}
}

\date{May 2026}

\begin{document}

\maketitle

\begin{abstract}
Current AI testing paradigms ask \emph{did the function pass?}  We argue that for
autonomous systems the correct question is: \emph{is this capability real, repeatable,
recoverable, and certifiable?}  We formalise this intuition with the
\textbf{Capability-Driven Intelligence Certification} (CDIC) framework, a six-stage
lifecycle model (DEFINED $\to$ TESTED $\to$ CERTIFIED $\to$ MONITORED $\to$ DEGRADED
$\to$ REVOKED) in which trust in a capability is not permanent but is continuously
earned, monitored, and can be revoked.  Each capability carries a score $s \in [0,1]$
that must exceed the sovereign deployment threshold $\varphi^{-1} \approx 0.618$ to
remain certified.  The framework introduces five Domain-Specific Testing Languages
(CTL, MTL, WTL, ATL, ETL) aligned to the five NOVA CPL-F substrate layers, and defines
proof-linked certificates that embed test logs, runtime evidence, and revocation rules.
We prove that the CDIC lifecycle is a monotone Markov chain on the lattice of trust
states, and that self-healing behaviour corresponds to upward transitions in this lattice.
Empirical results from the NOVA Alpha Test Suite (1734 assertions, 100\% pass rate under
current test conditions) populate the SVA evidence matrix.  We classify outstanding
claims and outline the experimental programme needed to extend verified results to
emergent-behaviour domains.
\end{abstract}

\tableofcontents
\newpage

% ─────────────────────────────────────────────────────────────
\section{Introduction}
\label{sec:intro}
% ─────────────────────────────────────────────────────────────

Autonomous AI systems present a certification challenge that traditional software
testing does not address.  A function that passes its unit tests today may fail under
novel inputs tomorrow; a self-modifying agent that learns may acquire unanticipated
behaviours; a fleet of cooperating agents may exhibit emergent failure modes that no
individual agent displays alone.

The NOVA sovereign AGI architecture \cite{nova_paper3} addresses this challenge through
a continuous, revocable trust model.  Rather than certifying an agent once and treating
certification as permanent, the NOVA model certifies individual \emph{capabilities}—
well-defined behavioural properties—and monitors them continuously.  A capability may
be \emph{discovered} through testing, \emph{certified} when evidence meets threshold,
\emph{degraded} when runtime monitoring detects anomaly, and \emph{revoked} when the
capability can no longer be demonstrated.

This paper makes four contributions:
\begin{enumerate}
  \item The \textbf{CDIC lifecycle model} (§2): a six-state trust automaton with
    $\varphi$-encoded score thresholds.
  \item \textbf{Domain-Specific Testing Languages} (§3): CTL, MTL, WTL, ATL, and ETL,
    one per NOVA substrate layer.
  \item \textbf{Formal properties} (§4): the CDIC lattice is a monotone Markov chain;
    self-healing is upward lattice movement.
  \item \textbf{Claim classification and evidence matrix} (§5): mapping all major
    NOVA testing assertions to their appropriate epistemic status.
\end{enumerate}

% ─────────────────────────────────────────────────────────────
\section{The CDIC Lifecycle Model}
\label{sec:lifecycle}
% ─────────────────────────────────────────────────────────────

\subsection{Capability Definition}

\begin{definition}[Capability]
A \emph{capability} $C = (I, P, S, T)$ is a tuple where:
\begin{itemize}[noitemsep]
  \item $I$ is a unique identifier (format: \texttt{CAP-NAME-NNN}),
  \item $P$ is a predicate over observable system states,
  \item $S \in [0,1]$ is the current score derived from test evidence,
  \item $T \in \mathcal{L}$ is a trust state in the lifecycle lattice $\mathcal{L}$.
\end{itemize}
\end{definition}

\begin{definition}[Lifecycle Lattice]
The \emph{capability lifecycle lattice} is the totally ordered set
\[
  \mathcal{L} = (\text{DEFINED} \prec \text{TESTED} \prec
                 \text{CERTIFIED} \prec \text{MONITORED} \prec
                 \text{DEGRADED} \prec \text{REVOKED}),
\]
with partial order $\prec$ and bottom element DEFINED, top element REVOKED.
\end{definition}

\begin{remark}
REVOKED is the top element because it represents the highest certainty that a
capability is \emph{absent}.  The ordering captures the epistemic trajectory of trust:
from \emph{undefined} through \emph{demonstrated} to \emph{certified} (peak trust)
and back down through \emph{monitored concern} and \emph{degradation} to \emph{revocation}.
\end{remark}

\subsection{Score Thresholds}

The sovereign deployment threshold is $\tau_\varphi = \varphi^{-1} \approx 0.6180$.
A capability may be certified only if $S \ge \tau_\varphi$.

\begin{definition}[Certification Predicate]
$\mathsf{certifiable}(C) \iff S(C) \ge \varphi^{-1} \wedge T(C) \in \{\text{TESTED},\text{MONITORED}\}$.
\end{definition}

\begin{definition}[Revocation Predicate]
$\mathsf{revocable}(C) \iff S(C) < \varphi^{-2} \vee \neg P(\text{recent evidence})$.
\end{definition}

The choice of $\varphi^{-1}$ as the certification floor and $\varphi^{-2}$ as the
revocation floor creates a \emph{hysteresis band} of width $\varphi^{-1} - \varphi^{-2}
= \varphi^{-1}(1 - \varphi^{-1}) = \varphi^{-1} \cdot \varphi^{-2} = \varphi^{-3}$
(using $1 - \varphi^{-1} = \varphi^{-2}$).

\subsection{Lifecycle Transitions}

\begin{figure}[h]
\centering
\begin{tabular}{c}
DEFINED $\xrightarrow{\text{first test}}$ TESTED $\xrightarrow{S \ge \tau_\varphi}$
CERTIFIED $\xrightarrow{\text{monitoring on}}$ MONITORED \\[6pt]
MONITORED $\xrightarrow{S \text{ drops}}$ DEGRADED $\xrightarrow{S < \varphi^{-2}}$
REVOKED \\[6pt]
DEGRADED $\xrightarrow{\text{self-heal}}$ MONITORED $\xrightarrow{\text{re-cert}}$
CERTIFIED
\end{tabular}
\caption{CDIC lifecycle transitions.  Self-healing is the upward transition DEGRADED
$\to$ MONITORED $\to$ CERTIFIED.}
\end{figure}

\subsection{Certificate Structure}

\begin{definition}[Capability Certificate]
A \emph{certificate} for capability $C$ is a record:
\[
  \text{CERT}(C) = \bigl(I, S, T, \text{ProofRef}, \text{IssuedAt}, \text{ExpiresAt},
                          \text{RevocationRule}\bigr),
\]
where $\text{ProofRef}$ is a content-addressed reference to the test log and
$\text{RevocationRule}$ is a predicate that, if satisfied, triggers automatic
downgrade to REVOKED.
\end{definition}

% ─────────────────────────────────────────────────────────────
\section{Domain-Specific Testing Languages}
\label{sec:dsls}
% ─────────────────────────────────────────────────────────────

Each NOVA substrate layer has a corresponding testing language:

\begin{table}[h]
\centering
\caption{CDIC Domain-Specific Testing Languages}
\label{tab:dsls}
\begin{tabular}{lllp{5cm}}
\toprule
DSL & Name & Substrate Layer & Primary Test Types \\
\midrule
CTL & Capability Testing Language &
  CPL-F Math / Motoko &
  φ-constant verification, protocol invariant checking,
  canister state assertions \\
MTL & Monte Carlo Testing Language &
  CPL-F Workers / Substrate &
  Statistical convergence, walk energy, bootstrap CIs,
  stochastic stability bounds \\
WTL & Worker Testing Language &
  CPL-F SERVITORES Fleet &
  COR\_PARVUM timing (873ms), CEREBRUM coherence,
  MACHINA\_VIRTUALIS state transitions \\
ATL & Artifact Testing Language &
  CPL-F Frontend / Protocol Artifacts &
  File integrity, signature verification,
  artifact payload consistency, HTML volume \\
ETL & Endurance Testing Language &
  Cross-Layer Temporal &
  Long-term Fibonacci stability, Kuramoto drift,
  HEARTBEAT consistency, memory persistence,
  mathematical constant precision \\
\bottomrule
\end{tabular}
\end{table}

\subsection{CTL: Capability Testing Language}

CTL tests verify that a claimed capability is \emph{present} at the system level.
An assertion in CTL has the form:
\[
  \mathsf{assert\_cap}(I, P, \tau) \stackrel{\text{def}}{=}
  \text{``run test suite for capability } I,
  \text{ verify predicate } P, \text{ require score} \ge \tau\text{''}
\]

\subsection{MTL: Monte Carlo Testing Language}

MTL tests apply statistical testing to sovereign constants.  An MTL assertion:
\[
  \mathsf{assert\_mc}(f, n, \varepsilon) \stackrel{\text{def}}{=}
  |\hat{f}_n - f_\infty| < \varepsilon
  \quad \text{with probability} \ge 1 - \alpha
\]
where $f_\infty$ is the true value, $\hat{f}_n$ is the Monte Carlo estimate at $n$
samples, and $\varepsilon$ is the tolerance.

\subsection{WTL: Worker Testing Language}

WTL tests verify SERVITORES fleet autonomy.  Each worker $w$ with kernel ID
\texttt{GOL-XXX-NNN} must satisfy:
\begin{align*}
  &\mathsf{HB}(w) = 873 \text{ ms} \pm 5\text{ ms (COR\_PARVUM timing)}, \\
  &\mathsf{KID}(w) \text{ matches } /^\mathtt{GOL-[A-Z]+-\d{3}}$/, \\
  &\mathsf{FAMILY}(w) \in \text{LATIN\_FAMILY\_NAMES}.
\end{align*}

\subsection{ATL: Artifact Testing Language}

ATL tests verify that deployed artifacts are consistent with declared specifications:
\[
  \mathsf{assert\_artifact}(a, \text{type}, \text{size\_lb}) \stackrel{\text{def}}{=}
  |a| \ge \text{size\_lb} \wedge \mathsf{type}(a) = \text{type}
\]

\subsection{ETL: Endurance Testing Language}

ETL tests verify temporal stability.  An ETL assertion specifies a sequence of $n$
ticks and a stability predicate $Q$:
\[
  \mathsf{assert\_endure}(Q, n) \stackrel{\text{def}}{=}
  \bigwedge_{t=1}^{n} Q(\text{state}_t)
\]

% ─────────────────────────────────────────────────────────────
\section{Formal Properties of CDIC}
\label{sec:formal}
% ─────────────────────────────────────────────────────────────

\subsection{Monotone Markov Chain}

\begin{definition}[CDIC Markov Chain]
The CDIC process is a Markov chain $(T_n)_{n \ge 0}$ on $\mathcal{L}$ with transition
probabilities $p_{ij} = P(T_{n+1} = j \mid T_n = i)$.
\end{definition}

\begin{theorem}[Monotonicity]
\label{thm:mono}
The CDIC Markov chain is stochastically monotone: if $i \prec j$ in $\mathcal{L}$,
then the distribution of $T_{n+1}$ starting from state $j$ first-order stochastically
dominates the distribution starting from state $i$.
\end{theorem}

\begin{proof}[Proof Sketch]
Score $S$ decreases monotonically under no-maintenance; upward transitions require
deliberate remediation.  The transition matrix is upper-triangular on the lattice
(capabilities can only improve with intervention or worsen naturally), giving
stochastic monotonicity via standard results \cite{lindvall1992lectures}.
\end{proof}

\subsection{Self-Healing as Lattice Recovery}

\begin{definition}[Self-Healing Event]
A \emph{self-healing event} is a transition $T_n \to T_{n+1}$ with
$T_{n+1} \prec T_n$ in $\mathcal{L}$ (downward in the degradation direction).
\end{definition}

\begin{theorem}[Self-Healing Condition]
\label{thm:heal}
A capability $C$ with $T(C) = \text{DEGRADED}$ can undergo self-healing if and only if
its self-healing mechanism $H_C$ produces evidence $e$ satisfying $P(e) = \text{true}$
and the resulting score $S' \ge \tau_\varphi$.
\end{theorem}

\subsection{Capability Independence}

\begin{theorem}[Certificate Independence]
Two capabilities $C_1, C_2$ with disjoint predicate domains ($\text{dom}(P_1) \cap
\text{dom}(P_2) = \emptyset$) have independent lifecycle trajectories.  That is,
$T_n^{(1)}$ and $T_n^{(2)}$ are independent processes.
\end{theorem}

This justifies the modular certification architecture: each capability is certified
independently, and revocation of one does not automatically imply revocation of others.

\subsection{$\varphi$-Encoded Hysteresis}

\begin{proposition}
The hysteresis band $[\varphi^{-2}, \varphi^{-1}]$ has width $\varphi^{-3}$, the third
power of $\varphi^{-1}$.
\end{proposition}

\begin{proof}
$\varphi^{-1} - \varphi^{-2} = \varphi^{-1}(1 - \varphi^{-1}) = \varphi^{-1} \cdot \varphi^{-2}
= \varphi^{-3}$, using the identity $1 - \varphi^{-1} = \varphi^{-2}$.
\end{proof}

The use of consecutive powers of $\varphi^{-1}$ for certification ($\varphi^{-1}$) and
revocation ($\varphi^{-2}$) thresholds ensures that the hysteresis width is itself
$\varphi$-encoded, creating a self-similar trust boundary.

% ─────────────────────────────────────────────────────────────
\section{Claim Classification and Evidence Matrix}
\label{sec:evidence}
% ─────────────────────────────────────────────────────────────

\subsection{Classification Scheme}

\begin{enumerate}[label=\textbf{\arabic*.}]
  \item \textbf{Verified (V)}: Directly observable from test logs or code.
  \item \textbf{Supported (S)}: Backed by experimental results under stated conditions.
  \item \textbf{Hypothesis (H)}: Plausible conjecture requiring comparative baselines.
  \item \textbf{Thesis (T)}: Architectural or strategic framing claim.
\end{enumerate}

\subsection{Evidence Matrix}

\begin{table}[h]
\centering
\caption{NOVA capability evidence matrix (selected entries)}
\label{tab:evidence}
\begin{tabular}{p{5.5cm}lll}
\toprule
Claim & Class & Test Suite Ref & Status \\
\midrule
PHI constants equal expected values & V & §21.2, §23.5, §1 & CERTIFIED \\
Parsers pass test suites & V & §7--§9 & CERTIFIED \\
Stress framework tests pass & V & §17 & CERTIFIED \\
Worker family registry correct & V & §16, §20 & CERTIFIED \\
DSL parsers exist & V & §19 & CERTIFIED \\
1734 tests pass at 100\% & S & All §§ & CERTIFIED \\
$\pi$ convergence $O(n^{-1/2})$ & S & §21.1 & CERTIFIED \\
AR(1) AMOR-damped stability & S & §21.4 & CERTIFIED \\
Bootstrap CI shrinkage & S & §21.5 & CERTIFIED \\
Lyapunov guard discriminates chaos & S & §22.5 & CERTIFIED \\
Memory φ-decay persistent $>24h$ & S & §23.4 & CERTIFIED \\
HEARTBEAT regularity $\pm 8$ ms & S & §23.3 & CERTIFIED \\
$\varphi$ universally optimal constant & H & None yet & TESTED \\
$\varphi$-optimality better than $e$ & H & Comparative missing & DEFINED \\
Self-healing 99.7\% under all faults & H & Limited fault set & TESTED \\
Emergent behaviour > 99\% covered & H & No emergent baseline & DEFINED \\
Testing as immune system metaphor & T & N/A (framing) & N/A \\
Capabilities $>$ functions & T & N/A (architectural) & N/A \\
\bottomrule
\end{tabular}
\end{table}

\subsection{Deployment Readiness Rules}

A component is \emph{deployment-ready under supervised conditions} if:
\begin{enumerate}[noitemsep]
  \item All capabilities relevant to its function are in state CERTIFIED or MONITORED.
  \item Mean capability score $\bar{S} \ge \varphi^{-1}$.
  \item No capability has state REVOKED or DEGRADED.
  \item Self-healing mechanisms are tested and carry state $\ge$ TESTED.
\end{enumerate}

\begin{remark}[External Release Boundary]
The following language is appropriate for external release of deployment claims:

\emph{``In our current implementation, the validation suite reports 1734 passing tests
across language parsers, Monte Carlo simulation, stress testing, backend engines, model
routing, memory endurance, and AI capability checks.  These results support supervised
deployment readiness for the tested components.  Broader claims around emergent
behaviour, $\varphi$-universality, and production autonomy require continued external
benchmarking.''}
\end{remark}

% ─────────────────────────────────────────────────────────────
\section{Integration with the Sovereign Validation Authority}
\label{sec:sva_integration}
% ─────────────────────────────────────────────────────────────

The CDIC framework is operationalised by the Sovereign Validation Authority (SVA)
\cite{nova_sva}, which:

\begin{itemize}[noitemsep]
  \item Maintains the capability registry and certificate database.
  \item Runs the five DSL test suites (CTL, MTL, WTL, ATL, ETL) on every build.
  \item Applies the lifecycle transitions of §2.3 based on test results.
  \item Enforces the deployment readiness rules of §5.3.
  \item Provides continuous monitoring via integration with the CPL/PULSE runtime.
\end{itemize}

The SVA embodies the core CDIC doctrine:

\begin{center}
\emph{No capability is real until it is tested, proof-linked, monitored, and revocable.}
\end{center}

% ─────────────────────────────────────────────────────────────
\section{Related Work}
\label{sec:related}
% ─────────────────────────────────────────────────────────────

Testing frameworks for AI systems have been proposed in the context of
property-based testing \cite{quickcheck}, runtime monitoring \cite{leucker2009brief},
and formal verification \cite{katz2017reluplex}.  However, these approaches treat
capabilities as binary (pass/fail) and do not support the lifecycle model of
trust accumulation and revocation that autonomous systems require.

The closest antecedent is the \emph{runtime assurance} literature \cite{schierman2015},
which inserts monitoring contracts into safety-critical systems.  CDIC extends this to
include explicit score dynamics, $\varphi$-encoded thresholds, and multi-level
revocability.

% ─────────────────────────────────────────────────────────────
\section{Conclusion}
\label{sec:conclusion}
% ─────────────────────────────────────────────────────────────

We have presented the Capability-Driven Intelligence Certification (CDIC) framework,
formalising the lifecycle of trust in autonomous AI capabilities as a monotone Markov
chain on a $\varphi$-encoded trust lattice.  Key contributions include:

\begin{itemize}
  \item A six-state lifecycle automaton with $\varphi^{-1}$ certification and
    $\varphi^{-2}$ revocation thresholds, creating a self-similar hysteresis band.
  \item Five domain-specific testing languages (CTL, MTL, WTL, ATL, ETL) aligned to
    the NOVA CPL-F substrate layers.
  \item Formal proofs that CDIC is stochastically monotone and that self-healing
    corresponds to downward lattice movement.
  \item A claim classification scheme (V/S/H/T) with an evidence matrix linking
    all major NOVA capability assertions to test evidence.
\end{itemize}

Outstanding research includes the comparative baseline experiments needed to elevate
$\varphi$-universality from hypothesis (H) to supported (S), and the design of
emergent-behaviour test protocols.

\begin{thebibliography}{9}
\bibitem{nova_paper3}
  NOVA Sovereign Research,
  ``Self-Healing Multi-Agent Systems,''
  \texttt{docs/papers/arxiv/paper3\_self\_healing\_multi\_agent\_systems.tex}, 2024.

\bibitem{lindvall1992lectures}
  T.\ Lindvall,
  \emph{Lectures on the Coupling Method},
  Wiley, 1992.

\bibitem{quickcheck}
  K.\ Claessen, J.\ Hughes,
  ``QuickCheck: A Lightweight Tool for Random Testing of Haskell Programs,''
  \emph{ICFP}, 2000.

\bibitem{leucker2009brief}
  M.\ Leucker, C.\ Schallhart,
  ``A Brief Account of Runtime Verification,''
  \emph{Journal of Logic and Algebraic Programming}, 78(5):293--303, 2009.

\bibitem{katz2017reluplex}
  G.\ Katz et al.,
  ``Reluplex: An Efficient SMT Solver for Verifying Deep Neural Networks,''
  \emph{CAV}, 2017.

\bibitem{schierman2015}
  J.\ Schierman et al.,
  ``Runtime Assurance for Autonomous Systems,''
  \emph{AIAA Infotech}, 2015.

\bibitem{nova_paper7}
  NOVA Sovereign Research,
  ``Kuramoto AGI Reasoning,''
  \texttt{docs/papers/arxiv/paper7\_kuramoto\_agi\_reasoning.tex}, 2025.

\bibitem{nova_sva}
  NOVA Sovereign Research,
  ``Sovereign Validation Authority Charter v1,''
  \texttt{docs/charters/SVA\_CHARTER.md}, 2026.

\bibitem{nova_paper10}
  NOVA Sovereign Research,
  ``Monte Carlo Verification of $\varphi$-Optimality in Sovereign AGI Systems,''
  \texttt{docs/papers/arxiv/paper10\_monte\_carlo\_phi\_optimality.tex}, 2026.
\end{thebibliography}

\end{document}
